\id{ҒТАМР 06.39.02}{}

\begin{header}
\swa{З.Т.~Төребекова, Б.Н.~Боқаев, А.Д.~Бактияров, Н.Б.~Айжарыков}{ЦИФРЛАНДЫРУ ЖӘНЕ ЖАСАНДЫ ИНТЕЛЛЕКТ ҚАЗАҚСТАНДАҒЫ ЖОҒАРЫ ОҚУ ЖЕРЛЕРІНІҢ ЖАҺАНДЫҚ БӘСЕКЕЛЕСТІККЕ ҚАБІЛЕТТІЛІГІН АРТТЫРУ ФАКТОРЫ РЕТІНДЕ: ЭКОНОМИКАЛЫҚ ЖӘНЕ ИНСТИТУЦИОНАЛДЫҚ АСПЕКТІЛЕР}

\tsp{1}З.Т.~Төребекова\envelope,
\tsp{2}Б.Н.~Боқаев,
\tsp{3}А.Д.~Бактияров,
\tsp{1}Н.Б.~Айжарыков
\end{header}

\begin{affil}
\tsp{1}Алматы менеджмент университеті, Алматы, Қазақстан,

\tsp{2}Қарағанды индустриалды университеті, Теміртау, Қазақстан,

\tsp{3}Еуразия ұлттық университеті, Астана, Қазақстан

\corrauthor{Корреспондент-автор: zttemirkhan@gmail.com}
\end{affil}

Зерттеу тақырыбының өзектілігі цифрлық технологиялар мен жасанды
интеллектінің (ЖИ) қар\-қынды дамуы жоғары оқу орындарындағы білім беру
және басқарушылық үдерістерді түбегейлі өзгертетіндігінде, бұл
университеттердің жаһандық бәсекеге қабілеттілігін арттыруға мүмкіндік
беретінімен айқындалады. Зерттеудің мақсаты -- Қазақстанның жоғары және
жоғары оқу орнынан кейінгі білім беру жүйесіндегі цифрлық
трансформацияның жай-күйін зерттеу және жасанды интеллектіні енгізудің
ерекшеліктерін бағалау.

Осы мақсатқа қол жеткізу үшін мынадай міндеттер жүзеге асырылды:
университеттердің цифрлық жетілу деңгейін анықтау, нормативтік-құқықтық
базаны, цифрландырудың жоғары оқу орындарының экономикалық тұрақтылығына
және олардың нарықтық позицияларына әсерін талдау ондай-ақ цифрландыру
жолындағы негізгі кедергілер мен мүмкіндіктерді қарастыру. Зерттеу
объектісі ретінде Қазақстанның жоғары оқу орындарының қызметі алынып,
зерттеу пәні -- олардың цифрландыруға және жасанды интеллектіні
интеграциялауға бағытталған тәсілдері болып табылады. Зерттеу әдіснамасы
стратегиялық құжаттарды контент-талдау, ресми статистикалық деректерді
сандық өңдеу және салыстырмалы талдау әдістеріне негізделген.

Зерттеу нәтижелері Қазақстан университеттерінің цифрлық жетілу
деңгейінің қазіргі уақытта «орташа» деңгейде екенін көрсетті. Жасанды
интеллектіні енгізу операциялық шығындарды оңтайландыруға және білім
беру қызметтерінің экспорттық әлеуетін арттыруға ықпал ететіні
анықталды. Цифрлық технологияларды енгізудің ең тиімді тәсілі
-- технологиялық инфрақұрылымды дамытумен қатар академиялық ортадағы
цифрлық құзыреттіліктерді жүйелі түрде арттыруға бағытталған кешенді
трансформация әдісі болып табылады. Университеттердің қызметін
жетілдіруге арналған ұсыныстар ретінде оқытушылардың жасанды интеллект
құралдарымен жұмыс істеу дағдыларын дамытуға бағытталған кәсіби қайта
даярлау бағдарламаларын енгізу, IT-инфраструктураны қаржыландыру үшін
мемлекеттік-жекеменшік әріптестік тетіктерін әзірлеу және инновациялық
экожүйені институционалдық қолдауды қамтамасыз ету қажеттілігі
негізделген.

{\bfseries Түйін сөздер:} жоғары және жоғары оқу орнынан кейінгі білім беру
жүйесі, цифрландыру, жасанды интеллект, Қазақстан, жаһандық бәсекеге
қабілеттілік.

\begin{header}
ЦИФРОВИЗАЦИЯ И ИСКУССТВЕННЫЙ ИНТЕЛЛЕКТ КАК ФАКТОР ПОВЫШЕНИЯ ГЛОБАЛЬНОЙ КОНКУРЕНТОСПОСОБНОСТИ ВЫСШИХ УЧЕБНЫХ ЗАВЕДЕНИЙ КАЗАХСТАНА: ЭКОНОМИЧЕСКИЕ И ИНСТИТУЦИОНАЛЬНЫЕ АСПЕКТЫ

\tsp{1}З.Т.~Торебекова\envelope,
\tsp{2}Б.Н.~Бокаев,
\tsp{3}А.Д.~Бактияров,
\tsp{1}Н.Б.~Айжарыков
\end{header}

\begin{affil}
\tsp{1}Алматы менеджмент университет, Алматы, Казахстан,

\tsp{2} Карагандинский индустриальный университет, Темиртау, Казахстан,

\tsp{3} Евразийский национальный университет им. Л.Н. Гумилева, Астана, Казахстан,

e-mail: zttemirkhan@gmail.com
\end{affil}

Актуальность темы исследования определяется тем, что стремительное
развитие цифровых технологий и искусственного интеллекта (ИИ) радикально
меняет образовательные и управленческие процессы в высших учебных
заведениях, создавая возможности для повышения их глобальной
конкурентоспособности. Цель исследования -- изучить состояние цифровой
трансформации в системе высшего и послевузовского образования Казахстана
и оценить особенности внедрения искусственного интеллекта.

Для достижения данной цели были реализованы следующие задачи:
определение уровня цифровой зрелости университетов, анализ
нормативно-правовой базы, оценка влияния цифровизации на экономическую
устойчивость вузов и их рыночные позиции, а также рассмотрение ключевых
барьеров и возможностей на пути цифровой трансформации. Объектом
исследования является деятельность высших учебных заведений Казахстана,
а предметом исследования -- подходы к их цифровизации и интеграции
искусственного интеллекта. Методология исследования основана на
применении методов контент-анализа стратегических документов,
количественной обработки официальной статистики и сравнительного
анализа.

Результаты исследования показали, что уровень цифровой зрелости
казахстанских университетов в настоящее время находится на среднем
уровне. Выявлено, что внедрение искусственного интеллекта способствует
оптимизации операционных затрат и повышению экспортного потенциала
образовательных услуг. Наиболее эффективным подходом к внедрению
цифровых технологий является метод комплексной трансформации,
направленный на системное повышение цифровых компетенций в академической
среде параллельно с развитием технологической инфраструктуры. В качестве
рекомендаций по совершенствованию деятельности университетов обоснована
необходимость внедрения программ профессиональной переподготовки
преподавателей по работе с инструментами ИИ, разработки механизмов
государственно-частного партнерства для финансирования IT-инфраструктуры
и обеспечения институциональной поддержки инновационной экосистемы.

{\bfseries Ключевые слова:} система высшего и послевузовского образования,
цифровизация, искусственный интеллект, Казахстан, глобальная
конкурентоспособность.

\begin{header}
DIGITALIZATION AND ARTIFICIAL INTELLIGENCE AS A FACTOR IN ENHANCING THE GLOBAL COMPETITIVENESS OF HIGHER EDUCATION INSTITUTIONS IN KAZAKHSTAN: ECONOMIC AND INSTITUTIONAL ASPECTS

\tsp{1}Z.T.~Torebekova\envelope,
\tsp{2}B.N.~Bokaev,
\tsp{3}A.D.~Baktiyarov,
\tsp{1}N.B.~Aizharykov
\end{header}

\begin{affil}
\tsp{1}Almaty Management University, Almaty, Kazakhstan,

\tsp{2}Karaganda Industrial University, Temirtau, Kazakhstan,

\tsp{3}L.N. Gumilyov Eurasian National University, Astana, Kazakhstan,

e-mail: zttemirkhan@gmail.com
\end{affil}

The relevance of this study is determined by the rapid development of
digital technologies and artificial intelligence (AI), which
fundamentally transform the educational and administrative processes of
higher education institutions, thereby enhancing their global
competitiveness. The aim of the research is to examine the current state
of digital transformation in Kazakhstan's higher and postgraduate
education system and to assess the features of AI implementation.

To achieve this goal, the following objectives were implemented:
determining the level of digital maturity of universities, analyzing the
regulatory and legal framework, assessing the impact of digitaliza\-tion
on the economic sustainability and market positions of higher education
institutions, as well as considering key barriers and opportunities on
the path to digital transformation. The research object is the
activities of higher education institutions in Kazakhstan, while the
research subject is the approaches to their digitalization and
integration of artificial intelligence. The research methodology is
based on content analysis of strategic documents, quantitative
processing of official statistical data, and comparative analy\-sis
methods.

The results indicate that the digital maturity of Kazakhstan's
universities is currently at an ``average'' level. It has been
established that the integration of artificial intelligence contributes
to the optimization of operational costs and the enhancement of the
export potential of educational services. The most effective approach to
implementing digital technologies is identified as the method of
comprehensive transformation, which focuses on the systematic
advancement of digital competencies within the academic environment
alongside the development of technological infrastructure. As
recommendations for improving university operations, the study
substantiates the necessity of implementing professional retraining
programs for faculty in AI tools, developing public-private partnership
mechanisms for financing IT infrastructure, and ensuring institutional
support for the innovation ecosystem.

{\bfseries Keywords:} higher and postgraduate education system,
digitalization; artificial intelligence, Kazakhstan, global
competitiveness.

\begin{multicols}{2}
{\bfseries Кіріспе.} Цифрлық технологиялардың қарқынды дамуына және білім
экономикасының пайда болуына орай жаһандық ЖЖОКБ жүйелері ауқымды
құрылымдық және мазмұндық өзгерістерге ұшырауда. Осы мәселенің
өзектілігі университеттердің халықаралық білім беру нарығындағы орны
олардың цифрлық әлеуеті мен жасанды интеллектіні (ЖИ) қолдану деңгейіне
тікелей байланысты болуымен түсіндіріледі. Интеллектуалды ресурстар үшін
жаһандық бәсекелестік жағдайында әлемдік жетекші университеттер цифрлық
шешімдерді жаһандық бәсекеге қабілеттілікті нығайтуға арналған
стратегиялық актив және жоғары оқу орнының материалдық емес активтерін
капиталдандырудың маңызды құралы ретінде қарастырады {[}1{]}.

Қазіргі жағдайда ЖЖОКБ-ның жаһандық бәсекеге қабілеттілігі олардың
технологиялық өзгерістерге тез бейімделу, жоғары сапалы білім беру
қызметтерін ұсыну және шекті өнімділігі жоғары, сұранысқа ие адами
капиталды қалыптастыру қабілетімен анықталады. Халықаралық рейтингтер
мен үлкен деректерді талдау университеттер арасындағы бәсекелестікті
күшейтіп, цифрлық жетілуді білім беру мекемелерінің қаржылық тұрақтылығы
мен инвестициялық тартымдылығының маңызды көрсеткіштеріне айналдыруда
{[}2{]}. Осыған байланысты ЖЖОКБ-ді цифрландыру мемлекеттік саясаттың
басым бағыты ретінде қарастырылуда. Осы мақсатта ұлттық цифрлық даму
стратегиялары әзірленіп, ЖИ-ді білім беру және ғылыми қызметке
интеграциялауды көздейтін құқықтық және институционалдық жағдайлар
жасалуда, сондай-ақ инновацияларды қолдау және цифрлық құзыреттіліктерді
дамыту тетіктері белгіленуде.

Бұл ретте технологиялық прогресті ынталандыру, білім беру үдерісіне
қатысушылардың құқықтарын қорғау және жасанды интеллектіні пайдаланудың
этикалық қағидаттарын сақтау арасындағы тепе-теңдікті қамтамасыз ететін
нормативтік-құқықтық базалар ерекше маңызға ие. Қазақстан Республикасы
үшін цифрландыру және ЖЖОКБ-ге ЖИ-ні енгізу ұзақ мерзімді ұлттық
басымдықтарды жүзеге асыру және жаһандық білім беру кеңістігіне
интеграциялану тұрғысынан стратегиялық маңызға ие {[}3{]}. Осыған
байланысты Қазақстан үкіметі адами капиталды, инновациялық экономиканы
және бәсекеге қабілетті білім беру жүйесін дамытуға бағытталған
мемлекеттік бағдарламалар мен стратегиялық құжаттарда бекітілген цифрлық
трансформация бағытын дәйекті түрде жүзеге асыруда.

Алайда цифрлық трансформацияның табысты болуы білім беру мекемелерінің
цифрлық жетілу деңгейіне, олардың инфрақұрылымының, адами ресурстарының
жағдайына және ЖИ негізіндегі шешімдерді енгізуге қажетті реттеуші және
институционалдық ортаның дайындығына тікелей байланысты. Стратегиялық
мақсаттар мен цифрлық бастамаларды іске асырудың практикалық тетіктері
арасындағы үйлесімділіктің жеткіліксіздігі бұл шаралардың тиімділігін
төмендетуі және университеттердің жаһандық бәсекеге қабілеттілігін
арттыру әлеуетін шектеуі мүмкін {[}4{]}.

Қазіргі кездегі ғылыми әдебиетте цифрландыру университеттердің
бизнес-процестерін оңтайландырудың және трансакциялық шығындарды
төмендетудің негізгі факторы ретінде қарастырылады {[}5{]}. Экономикалық
тұрғыдан цифрлық трансформация университеттердің құндылық
қалыптастырудың дәстүрлі моделін түбегейлі өзгертіп, оны икемді әрі
ауқымды етеді {[}6{]}. Талапкерлер жоғары технологиялармен қамтамасыз
етілген ортаны таңдай отырып, университеттің репутациялық капиталы мен
нарықтағы ұзақ мерзімді кірістілігінің артуына әсер етеді {[}1{]}. ЖИ
құралдарын енгізу білім беру қызметтерінің білім беру өнімін терең
даралауды және жоғары операциялық тиімділік деңгейін қамтамасыз етеді
{[}7{]}. Бұл тәсілді цифрлық жетілу тұжырымдамасымен толықтырады,
ұйымның өзгермелі цифрлық бәсекелестік ортаға бейімделу қабілетін атап
көрсетеді {[}2{]}.

Бұл ретте университеттерді таңдауда талапкерлер жоғары технологиялармен
қамтамасыз етілген білім беру үдерісіне басымдық береді {[}8{]}. Осыған
байланысты университет басшылары студенттердің тәжірибесін жақсартуды
цифрлық трансформацияның негізгі нәтижесі ретінде қарастыруы тиіс
{[}9{]}. Студенттердің қанағаттануы университеттің беделін және білім
беру нарығындағы тартымдылығын арттырады, сондай-ақ бұл бәсекелестік
артықшылықтың дамуына ықпал етеді {[}1{]}.

Зерттеулерге сәйкес цифрландырудың басты мүмкіндіктерінің бірі --
білімге қолжетімділікті кеңейту және институционалдық икемділікті
арттыру {[}10{]}. Сонымен қатар әкімшілік және білім беру үдерістерін
автоматтандыру институционалдық процестердің ашықтығы мен тиімділігін
қамтамасыз етеді {[}7{]}. Жалпы алғанда цифрландыру қызмет көрсету
сапасы мен пайдаланушы тәжірибесін жақсартып, университеттердің бәсекеге
қабілеттілігін арттырады {[}11{]}.

ЖИ-ге келер болсақ, қазіргі зерттеушілер ЖИ-дің университеттердің білім
беру қызметін кешенді түрде жетілдіру әлеуетіне ден қоюда {[}12,13{]}.
ЖИ білім беру деректерін талдаудан бастап, білім алу дербестігіне
дейінгі кең ауқымды міндеттерді іске асырады {[}14{]}. Алайда зерттеулер
ЖИ-дің академиялық жетістіктерге әсері контекстке байланысты екенін атай
келе, университеттердің бәсекеге қабілеттілігін арттыруға қабілеттілігін
ең маңызды стратегиялық ресурс деп айқындайды {[}15{]}.

Осы орайда цифрлық инфрақұрылымның жеткіліксіздігі және ақпараттық
жүйелердің фрагменттілігі цифрлық бастамалардың тиімділігіне кері әсер
етеді {[}16{]}. Демек цифрландыру мен ЖИ-ді енгізу өздігінен
университеттердің бәсекелестік артықшылықтарына кепілдік бермейді
{[}17{]}.

Қазақстандық зерттеушілер ЖЖОКБ-ны цифрландыруды олардың бәсекеге
қабілеттілігін арттырудың негізгі факторы ретінде қарастырады.
Университеттердің білім беру процестерін жаңғыртуға, қолжетімділігін
кеңейтуге және инновациялық әлеуетін нығайтуға цифрлық технологиялардың
ықпал ететіні көрсетіледі {[}18{]}. Сонымен қатар цифрландыруды қолдану
білім беру сапасына және университеттердің беделіне елеулі әсер ететіні
анықталған {[}19{]}.

Жасанды интеллектіні қолдану тұрғысынан оқуды дербестендіру мен
аналитикалық шешімдер арқылы жаңа оқыту талаптарына бейімделу қабілеті
артады {[}20{]}. ЖИ-ді ұтымды пайдалану оқытушылардың цифрлық
құзыреттіліктерін жүйелі түрде дамытуға және үйлестірілген стратегия
қалыптастыруға мүмкіндік береді, бұл университеттердің бәсекеге
қабілеттілігінің орнықты өсуіне оң әсер етеді {[}4{]}. Бұл өз кезегінде
университеттердің бәсекеге қабілеттілігінің тұрақты өсуіне оң ықпалын
тигізеді. Осылайша, шетелдік және қазақстандық зерттеулер цифрландыру
және ЖИ енгізу университеттердің бәсекеге қабілеттілігінің маңызды
құрамдас бөліктері екенін растайды.

Осы орайда, Қазақстанның ЖЖОКБ жүйесіне ЖИ мен цифрландыруды енгізуге
қолайлы құқықтық және институционалдық ортаның қалыптасуын кешенді
талдаудың өзектілігі артуда. Бұл мақаланың мақсаты - Қазақстанның жоғары
оқу орындарының жаһандық бәсекеге қабілеттілігін арттыру факторлары
ретіндегі цифрлық технологиялар мен ЖИ-ды пайдаланудың негізгі
шарттарын, шектеулерін және мүмкіндіктерін анықтау, сондай-ақ ЖЖОКБ
мемлекеттік саясат пен басқару тәжірибесін жетілдіру жөніндегі ғылыми
негізделген ұсыныстарды тұжырымдау болып табылады.

\emph{Зерттеу нысаны} - Қазақстанның жоғары оқу орындары жүйесі.
\emph{Зерттеу пәні} - жоғары білім беру саласындағы цифрландыру және
жасанды интеллект технологияларын енгізу процестері.

\emph{Зерттеудің мақсаты} - қазақстандық жоғары оқу орындарының жаһандық
бәсекеге қабілеттілігін арттыру факторлары ретіндегі цифрлық
технологиялар мен жасанды интеллектіні (ЖИ) пайдаланудың негізгі
шарттарын, шектеулерін және мүмкіндіктерін анықтау, сондай-ақ
мемлекеттік саясат пен басқару тәжірибесін жетілдіру жөніндегі ғылыми
негізделген ұсыныстарды тұжырымдау болып табылады.

{\bfseries Материалдар мен әдістер.} Зерттеу әдіснамасы Қазақстан
Республикасының жоғары және жоғары оқу орнынан кейінгі білім беру
(ЖЖОКБ) жүйесінде цифрландыру мен жасанды интеллектіні (ЖИ) енгізудің
құқықтық және институционалдық ортасын бағалауға бағытталды. Зерттеу
Қазақстанның жоғары білім беру саласындағы мемлекеттік саясаттың
мазмұнын және университеттердің жаһандық бәсекеге қабілеттілігін арттыру
шарттарын айқындауға мүмкіндік беретін сапалық және сандық дереккөздерге
негізделді.

Зерттеу гипотезасы ЖИ-ді жүйелі түрде интеграциялау жоғары оқу орны
ресурстарының аллокативті тиімділігін (ресурстарды тиімді бөлуді)
арттырып, оның жаһандық нарықтағы тұрақты бәсекелестік артықшылықтарының
орнықты өсуіне ықпал етеді деген болжамға негізделді.

Зерттеудің негізгі сұрағы Қазақстанның ЖЖОКБ жүйесінде цифрландыру мен
жасанды интеллектіні енгізудің қандай құқықтық және институционалдық
шарттары университеттердің жаһандық бәсекеге қабілеттілігін арттырады
деген мәселені қамтыды. Зерттеу гипотезасы университеттердің цифрлық
жетілу деңгейі мен жасанды интеллектіні жүйелі түрде интеграциялау
олардың институционалдық тиімділігін арттырып, халықаралық білім беру
кеңістігіндегі тартымдылығы мен бәсекеге қабілеттілігінің орнықты өсуіне
ықпал етеді деген болжамға негізделді.

Зерттеу материалы ретінде ЖЖОКБ саласын цифрландыру мен жасанды
интеллектіні дамытуға қатысты нормативтік-құқықтық актілер, мемлекеттік
бағдарламалар мен стратегиялық құжаттар, сондай-ақ уәкілетті мемлекеттік
органдардың ресми статистикалық деректері пайдаланылды. Талдауға жасанды
интеллектіні дамыту тұжырымдамалары, жоғары білім беру жүйесінің цифрлық
трансформациясына арналған стратегиялық құжаттар, сондай-ақ білім беру
үдерісінде жасанды интеллектіні қолданудың институционалдық және
этикалық аспектілерін айқындайтын материалдар енгізілді. Зерттеу
материалы сапалық тұрғыдан құжаттардың мазмұнын, ал сандық тұрғыдан
білім беру бағдарламаларына жасанды интеллектіні енгізу көрсеткіштерін
және цифрлық бастамалардың қамтылу ауқымын сипаттады.

Зерттеу бірнеше кезеңде жүзеге асырылды. Бірінші кезеңде цифрландыру мен
жасанды интеллектінің жоғары білім беру жүйесіндегі рөліне қатысты
ғылыми тұжырымдамалар зерделенді. Келесі кезеңде нормативтік және
стратегиялық құжаттарға жүйелі талдау жүргізіліп, институционалдық
алғышарттар айқындалды. Қорытынды кезеңде цифрлық инфрақұрылым мен білім
беру бастамаларына, соның ішінде AI-Sana акселерациялық платформасына
қатысты мәліметтер талданып, университеттердің білім беру қызметіне
жасанды интеллектіні интеграциялаудың ағымдағы деңгейі бағаланды.

Зерттеу барысында контент-талдау, жүйелік және салыстырмалы талдау
әдістері қолданылды. Алынған нәтижелер Қазақстанның ЖЖОКБ жүйесінде
цифрландыру мен жасанды интеллектіні енгізудің құқықтық және
институционалдық негіздерін жүйелеуге, сондай-ақ университеттердің
жаһандық бәсекеге қабілеттілігін арттыруға ықпал ететін негізгі
факторларды айқындауға мүмкіндік берді.

{\bfseries Нәтижелер және талқылау.} Соңғы жылдары Қазақстан
Республикасындағы ЖЖОКБ цифрландыру білім беру процесінде ЖИ-ді енгізуді
институционалдандыруға бағытталған жүйелі мемлекеттік саясат аясында
дамып келеді. Бұл бағыттағы маңызды қадамның бірі - 2024 жылы Қазақстан
Республикасы Ғылым және жоғары білім министрлігінің ЖЖОКБ-де жасанды
интеллектті қолданудың стандартын әзірлеуі болды. Осының негізінде
бірыңғай нормативтік және әдістемелік нұсқаулықтар белгіленді {[}21{]}.
Аталмыш стандартта ЖИ-ді қолданудың этикалық аспектілеріне, оның ішінде
академиялық тұтастыққа, алгоритмдік ашықтыққа және жеке деректерді
қорғауға ерекше мән беріледі. Осы тұрғыда университетаралық стандарт
цифрлық инновацияны ынталандыру, дербес оқу мен сыни ойлауды қоса
алғанда, ЖЖОКБ-ның негізгі қағидаттарын сақтауға қатысты теңгерімді
механизм қызметін атқарады.

Бұл ретте, көрсетілген шаралар Қазақстанда ЖИ технологияларын дамытудың
кең ауқымды мемлекеттік стратегиясымен тұспа тұс келгенін атап кету
қажет. Мысалы, Қазақстанда ЖИ дамытудың 2024--2029 жылдарға арналған
тұжырымдамасы бұл технологияны экономикалық өсудің негізгі драйвері және
ұлттық адами капиталды жаңғырту факторы ретінде анықтады {[}22{]}.
Сонымен қатар, елімізде ЖИ саласында қабылданып жатқан стратегиялық
шараларының маңызды нәтижелерінің қатарына ЖИ саласындағы білім беру
бағдарламаларының әзірленуін жатқызуға болады. Бүгінгі таңда
қазақстандық университеттер қолданбалы ЖИ, деректерді талдау,
кибернетика, медицинадағы жасанды интеллект, ақылды технологиялар және
пәнаралық инженерлік шешімдер сияқты салаларды қамтитын жиырмадан астам
мамандандырылған бағдарламаларды қолданысқа енгізген. ЖИ-ге қатысты
пәндер елдің көптеген университеттеріне интеграцияланған және бұл білім
беру бағдарламаларының цифрлық трансформациясының ауқымдылығын
дәлелдейді {[}23{]}.

Алайда, ЖЖОКБ-ді цифрландыру бірқатар құрылымдық қиындықтармен қатар
жүретінін назарда ұстау маңызды. Осы бағыттағы негізгі шектеулердің бірі
- оқытушылардың цифрлық және педагогикалық құзыреттіліктерінің
жеткіліксіздігі болып табылады. Бұның салдары білім берудегі ЖИ әлеуетін
толық жүзеге асыра алмаумен байланысты. Осыған орай, ЖИ-ді пайдалануға
бағытталған кәсіби даму бағдарламаларын кеңінен дамыту маңызды. Қазіргі
кезде жүргізілеп жатқан шаралардың басты мақсаты - университеттердің
цифрлық экономикада жұмыс істеуге институционалдық дайындығын арттыру
{[}24{]}.

Білім берудегі цифрландыру экожүйесін қалыптастырудың тағы бір маңызды
элементі - AI-Sana бағдарламасы сияқты инновациялық және жеделдету
бастамаларын іске асыру болып табылады. Мұндай жобалар білім беру
процесін ЖИ саласында жүргізіліп жатқан ғылыми зерттеулермен және
практикалық әзірлемелермен біріктіруге, студенттер мен жас
зерттеушілердің қолданбалы құзыреттіліктерін дамытуға ықпал етуді
көздейді {[}25{]}.

Сондай-ақ, цифрлық және ЖИ технологияларын тиімді қолдану жолында тиісті
цифрлық әлеуеттің қалыптасуы аса маңызды. Қазіргі кезде қазақстандық
ЖЖОКБ ұйымдарының (ЖЖОКБҰ) цифрлық жетілу деңгейін бағалау мақсатында
Ғылым және жоғары білім министрлігінің 2024 жылғы 14 тамыздағы №408
бұйрығымен Жоғары және/немесе жоғары оқу орнынан кейінгі білім беру
ұйымдарының цифрлық жетілу деңгейін анықтау жөніндегі әдістемелік
ұсыныстары бекітілген {[}3{]}. Бұл құжат цифрлық жетілуді бағалаудың
базалық, орта, жоғары және көшбасшылық сияқты 4 деңгейін мен бес негізгі
компонентін анықтайды: инфрақұрылым мен жабдықтау, цифрлық
платформаларды интеграциялау, білім беру процесінде жасанды интеллектті
пайдалану, оқытушылар мен студенттердің цифрлық құзыреттіліктері және
инновациялық білім беру ортасы.

Осы ретте, цифрлық жетілудің негізгі алғышарттарының қатарында тиісті
инфрақұрылым мен құрал-жабдықтардың болуы бар. Эмпирикалық зерттеулерге
сүйенсек, бүгінде еліміздегі мегаполистік және өңірлік университеттер
арасындағы цифрландыру деңгейінде теңгерімсіздік сақталып отыр. Зерттеу
нәтижелері өңірлік және орталық жоғары оқу орындары арасында цифрлық
алшақтық (digital divide) бар екенін көрсетті, бұл білім беру қызметтері
нарығындағы экономикалық пайданың біркелкі бөлінбеуіне және нарықтың
сегменттелуіне әкеледі. Компьютерлік парктің 95\%-ы интернетке
қосылғанымен, оқытушылардың 50\%-ында мамандандырылған цифрлық
құзыреттіліктер тапшылығы байқалады, бұл IT-инфраструктураға салынған
инвестициялардың қайтарымдылығын (ROI) тежейді {[}4{]}. Халықаралық
тәжірибе көрсеткендей, ЖИ-ді стратегиялық басқаруға енгізу әкімшілік
шығыстарды 15-20\%-ға азайтып, білім беру процесінің автономиясы мен
рентабельділігін арттырады {[}14{]}.

Осы саладағы басқа да зерттеулерге қатысқан студент-респонденттердің
26,5\%-ы отандық университеттердегі цифрлық инфрақұрылымды «орташа»,
16\% - «нашар» және 18,3\% - «орташадан жоғары»/«өте жақсы» деп
бағалаған. Студенттердің басым бөлігі (66,2\%) ЖИ құралдарын оқу
барысында тұрақты пайдаланатынын хабарлаған. Жалпы алғанда зерттеу
нәтижелері ЖЖОКБҰ-дағы цифрландыру мен ЖИ-ні енгізу деңгейін «орташа»
деп қорытындылаған {[}4{]}. ҚР Стратегиялық зерттеулер институтының
талдауында қазақстандық студенттердің 76\%-ға жуығы білім алу барысында
ЖИ-ге тұрақты жүгінетінін, сондай-ақ оқытушылардың (шамамен 51\%) ЖИ
құралдарын аналитика жасау және бағалау барысында қолданатынын
анықталған {[}26{]}.

Халықаралық зерттеулер, мысалы, Ұлыбритания жағдайында жүргізілген
зерттеулер ел аумағында ЖИ-ді белсенді қолданатын студенттердің үлесі
2024 жылғы 66\%-дан 2025 жылы 92\%-ға дейін өскенін көрсеткен {[}27{]}.
Бұған қоса, АҚШ-тың жетекші университеттері ЖИ құралдарын білім беру
үрдісін стратегиялық басқаруда белсенді пайдаланады {[}28{]}. Америкалық
элиталық колледж студенттерінің 80\%-дан астамы ЖИ-ді материалды
түсіндіру және күнделікті оқу тапсырмаларын орындау үшін ұдайы қолданады
{[}29{]}.

Цифрлық жетілудің тағы бір көрсеткіші -- цифрлық құзыреттіліктің
қалыптасуы болып табылады. Осы тұста, қазақстандық оқытушылардың цифрлық
құзыреттіліктерінің деңгейін бағалаған зерттеулердің қорытындысына
сәйкес, оқытушылардың 50\%-ға жуығы цифрлық және ЖИ құралдарын белсенді
пайдалануға қабілетті болса, қалған бөлігі тек базалық цифрлық
дағдыларға ие. Ғалымдардың пікірінше, бұл шектеулерді жою үшін
мамандандырылған кәсіби даму бағдарламалары және халықаралық
университеттермен тәжірибе алмасулар қажет. Қазіргі таңда цифрлық
құзыреттіліктерді дамыту негізінен ЖЖОКБҰ-ның басшыларын
университеттерді цифрлық трансформацияға дайындау бағдарламалары аясында
жүзеге асырылады {[}30{]}. Қазақстандық ЖЖОКБҰ-ғы цифрлық
құзыреттіліктің жеткіліксіздігі, ЖИ-дің бөлшектенген интеграциясының,
институционалдық қолдаудың кемшілігінен туындайтын салдар деп анықталған
{[}19{]}. Бұған қоса, оқытушылардың цифрлық құзыреттілігі LMS
жабдықтарының жеткіліксіздігі және ауқымды аналитикалық қызметтерге
қолжетімділіктің шектеулілігі себебінен дамымай отыр {[}4{]}.

Халықаралық тәжірибеде, мысалы Австралияның ЖЖОКБҰ-лары цифрлық
жетілудің жетекші деңгейіне стратегиялық жоспарлау, ЖИ-ді білім беру
және басқару процестеріне кешенді интеграциялау, тәуекелдерді басқару
және инновацияларды жүйелік қолдау арқылы қол жеткізген. ЖИ-ды жүйелі
пайдалану шетелдік оқу орындарында білім алу уақытын шамамен 27\%-ға
қысқартуға және оқу процесінің дербестігін арттыруға мүмкіндік берген
{[}1{]}. Сондай-ақ, шетелдік зерттеушілер ЖИ-ді пайдаланудың этикалық
тұстарын академиялық тұтастықпен теңдестірудің маңыздылығын да атап
көрсеткен {[}27{]}.

Осылайша, жүргізілген талдау қазақстандық ЖЖОКБҰ-дың цифрлық жетілуін
ішінара бағалауға мүкіндік берді. Мұндағы көрсеткіштер ретінде тиісті
зерттеулерде анықталған және ресми жарияланған деректер алынған.
Нәтижесінде, бүгінгі таңдағы қазақстандық университеттердің цифрлық
жетілу деңгейін орташа деп бағалауға болады.

Төмендегі 1-кестеде көрсетілгендей, қазақстандық университеттердің
цифрлық кемелдік деңгейі бес негізгі компонент бойынша «орташа» деп
бағаланады.
\end{multicols}

\tcap{1-кесте. Эмпирикалық деректер мен халықаралық көрсеткіштер
негізінде Қазақстан университеттерінің цифрлық жетілуін бағалау}
\begin{longtblr}[
  label = none,
  entry = none,
]{
  width = \linewidth,
  colspec = {Q[152]Q[206]Q[73]Q[312]Q[198]},
  cells = {c},
  cells = {font = \small},
  row{1} = {font=\bfseries\small},
  hlines,
  vlines,
}
Компонент & Қазақстан (эмпирикалық деректер) & Жетілу деңгейі & Экономикалық маңызы & Халықаралық бағдарлар\\
Инфрақұрылым және жабдықтар & 83 000 компьютер, оның 79 000 интернетке қосылған & Орта & Бұлтты есептеулер арқылы CAPEX пен OPEX-ті (күрделі және операциялық шығындарды) оңтайландыру & Жоғары өнімді есептеулерді (HPC) интеграциялау\\
Оқу процесінде жасанды интеллектіні қолдану & Студенттердің 76\%-ы ЖИ қолданады; оқытушылардың 51\%-ы & Орта & Дербестендіру – клиенттің (студенттің) өмірлік құндылығын (LTV) арттыру құралы ретінде & Бейімделген оқыту (Adaptive learning) жүйелерінің толық циклі\\
Оқытушылардың цифрлық құзыреттері & 50\% жуығы базалық цифрлық құзыреттерге ие & Орта & ЖИ-құзыреттіліктері болмаған жағдайда адами капиталдың құнсыздану (девальвация) тәуекелі & Үздіксіз даму (CPD) және ЖИ-сертификаттау\\
Инновациялық білім беру ортасы & Пилоттық жобалар, акселерациялық бағдарламалар (AI-Sana) & Орта & Зияткерлік меншікті (IP) коммерцияландыру & Технологиялық кәсіпкерлік экожүйесі
\end{longtblr}

\begin{multicols}{2}
Сонымен қатар, жүргізілген зерттеу нәтижелері қазіргі кезде қазақстандық
ЖЖОКБ жүйесіне цифрлық трансформацияны ұтымды іске асыруға қажетті
нормативтік-құқықтық базаның жеткілікті түрде қалыптасқанын көрсетті.
ЖИ-ды дамытудың ұлттық тұжырымдамасының және университетаралық
стандарттардың әзірленуі ЖИ кіріктірілген білім беру бағдарламаларының
ауқымын кеңейтуге және ЖЖОКБҰ-ны тұрақты цифрландыруға негіз береді.
Алайда, осы процестердің тиімділігі нормативтік талаптарды енгізу
сапасына, оқытушылар құрамының дайындық деңгейіне және этикалық
стандарттардың сақталуына тікелей байланысты.

{\bfseries Қорытынды.} Осы зерттеудің мақсаты Қазақстанның жоғары және
жоғары оқу орнынан кейінгі білім беру жүйесінде цифрлық технологиялар
мен жасанды интеллектіні (ЖИ) пайдаланудың университеттердің жаһандық
бәсекеге қабілеттілігіне әсерін анықтау болды. Зерттеу нәтижелері
қазақстандық университеттерде цифрлық трансформация процесі қалыптасу
және орташа деңгейдегі жетілу кезеңінде екенін көрсетті. Базалық цифрлық
инфрақұрылым болса да, цифрлық технологиялар мен ЖИ-ді институционалдық
дамудың стратегиялық ресурсы ретінде кешенді пайдалану жеткіліксіз
деңгейде қалып отыр {[}4{]}.

Сонымен қатар, жүргізілген зерттеу қазақстандық университеттердің
цифрлық трансформациясы ресурстарды экстенсивті жинақтау кезеңінде
екенін көрсетті. Бәсекеге қабілеттілікті нығайту үшін цифрландырудың
интенсивті моделіне көшу қажет. Бұл модель университеттердің
институционалдық орнықтылығын арттыруға, білімді тиімді монетизациялауға
және ЖЖОКБ жүйесін экономикалық трансформацияның локомотивіне
айналдыруға бағытталуы тиіс.

Эмпирикалық деректер университеттер арасында жасанды интеллектіге
негізделген шешімдерді енгізу деңгейінің біркелкі еместігін және цифрлық
құзыреттіліктердің дамуы бойынша асимметрияның бар екенін көрсетеді. Бұл
адами капиталды мақсатты түрде дамыту және цифрлық жетілуді басқаруда
жүйелі тәсіл қалыптастыру қажеттілігін айқындайды {[}30; 4{]}. Шетелдік
тәжірибе ЖИ технологияларының білім беру, ғылыми-зерттеу және әкімшілік
үдерістерге тікелей интеграциялануы университеттердің беделін арттыруға
оң әсерін тигізетінін растайды {[}27{]}.

Осыған байланысты Қазақстан университеттерінің жаһандық бәсекеге
қабілеттілігін арттыру үшін жекелеген цифрлық бастамалардан
институционалдық цифрлық трансформация моделіне көшу қажет. Бұл модель
оқытушылар мен әкімшілік персоналдың цифрлық құзыреттіліктерін жүйелі
дамытуға, басқарушылық шешімдер қабылдауда ЖИ негізіндегі аналитикалық
жүйелерді енгізуге, сондай-ақ цифрлық технологияларды білім сапасын
арттыру, ғылыми зерттеу өнімділігін жоғарылату және университет қызметін
халықаралық стандарттарға сәйкестендіру мақсатында пайдалануға
бағытталуы тиіс {[}4; 27{]}.

Зерттеу нәтижелері жоғары білім беру саласындағы мемлекеттік саясатты
жетілдіруде және университеттердің стратегиялық даму бағдарламаларын
әзірлеуде практикалық маңызға ие. Болашақ зерттеулер жасанды
интеллектіні пайдаланудың ұлттық стандарттарын әзірлеуге, этикалық және
құқықтық реттеуді жетілдіруге, сондай-ақ қазақстандық университеттердің
цифрлық жетілудің жоғары деңгейлеріне көшу жолдарын эмпирикалық тұрғыдан
зерделеуге бағытталуы мүмкін {[}13{]}.

\emph{{\bfseries Қаржыландыру.} Бұл зерттеу Қазақстан Республикасы Ғылым
және жоғары білім министрлігінің Ғылым комитеті тарапынан
қаржыландырылды (BR28712521).}
\end{multicols}

\begin{center}
{\bfseries Әдебиеттер}
\end{center}

\begin{refs}
1. George B., \& Wooden O. Managing the Strategic Transformation of
Higher Education through Artificial Intelligence//~Administrative
Sciences. -2023.-~Vol.13(9). -P.196. DOI 10.3390/admsci13090196.

2. Kane G. C. Digital maturity, not digital transformation// MIT Sloan
Management Review. -2017.- Vol.1(1). -P.1--15. URL:
\href{https://sloanreview.mit.edu/article/digital-maturity-not-digital-transformation}{https://sloanreview.mit.edu}.
- Date of access: 06.01.2026.

3. Zakon.kz. Қазақстан Республикасының Жоғары және (немесе) жоғары оқу
орнынан кейінгі білім беру ұйымдарының цифрлық жетілу деңгейін айқындау
жөніндегі әдістемелік ұсынымдарды бекіту туралы// Қазақстан Республикасы
Ғылым және жоғары білім министрінің міндетін атқарушының 2024 жылғы 14
тамыздағы № 408 бұйрығы. -2024. -
URL: \url{https://online.zakon.kz/Document/?doc_id=34078184}. -Қаралған
күні: 06.01.2026.

4. Джанегизова А., Huong Lan D., \& Чуланова З. Цифровая трансформация
высшего образования в Казахстане: проблемы и решения// Статистика, учет
и аудит\emph{.} - 2025. -Т.3(98). - С.249-263.DOI
10.51579/1563--415.2025.-3.18.

5. Antonopoulou K., Begkos C., \& Zhu Z. Staying afloat amidst extreme
uncertainty: A case study of digital transformation in higher
education// Technological Forecasting \& Social Change.
-2023.-Vol.192:122603. DOI 10.1016/j.techfore.2023.122603.

6. Hess T., Matt C., Benlian A., \& Wiesböck F. Options for formulating
a digital transformation strategy. MIS Quarterly Executive. -2016.-
Vol.15(2). -P.123-139. URL:
\url{https://aisel.aisnet.org/cgi/viewcontent.cgi?article=1335&context=misqe}.
Date of access: 06.01.2026.

7. Shrivastava S. K., \& Shrivastava C. New reforms in Indian higher
education: Rethinking the education model//Indian Journal of Social
Science and Literature. -2024.-Vol.4(1). -P.42-48.\\
DOI 10.54105/ijssl.E1142.04010924.

8. Yesner R. The Future of Higher Education: Digital Transformation Is
Critical to Learner and Institution Success// IDC.-2020.-URL:\\
\url{https://www.ecampusnews.com/files/2020/10/IDC_The-Future-of-Higher-Education.pdf}
-Date of access: 06.01 2026.

9. Stokes P., Baker N., DeMillo R., Soares L., \& Yaeger L. The
transformation-ready higher education institution. Huron; American
Council of Education; Georgia Institute of Technology. -2019.- URL:
\url{https://share.google/wbxeAYP6ZYMYlnZSn}.- Date of access: 06.01
2026.

10. Rodríguez D., Arias P., \& Jimenez M. Use of virtual resources for
the development of spatial vision in high school students//Cadernos de
Educação Tecnologia e Sociedade. -2023.- Vol.16(3). DOI
10.14571/brajets.v16.n3.473-489.

11. Yureva O. V., Burganova L. A., Kukushkina O. Y., Myagkov G. P., \&
Syradoev D. V. Digital transformation and its risks in higher education:
Students' and teachers' attitude//Universal Journal of Educational
Research.-2020.-Vol.8(11B).-P.5965-5971. DOI 10.13189/ujer.2020.082232.

12. Wang S., Wang F., Zhu Z., Wang J., Tran T., \& Du Z. Artificial
intelligence in education: A systematic literature review//Expert
Systems with Applications. -2024.- Vol.252:124167.\\
DOI 10.1016/j.eswa.2024.124167.

13. Bond M., Khosravi H., De Laat M., et al. A meta systematic review of
artificial intelligence in higher education: A call for increased
ethics, collaboration, and rigour//International Journal of Educational
Technology in Higher Education\emph{.} -2024.-Vol.21:4. DOI
10.1186/s41239-023-00436-z.

14. Boussouf Z., Amrani H., Zerhouni Khal M., \& Daidai F. Artificial
intelligence in education: A systematic literature review// Data and
Metadata\emph{.} -2024.- Vol.3: 288. DOI 10.56294/dm2024288.

15. Tlili A. Can artificial intelligence (AI) help in computer science
education? A meta-analysis approach// Revista Española de Pedagogía.
-2024. DOI 10.22550/2174-0909.4172.

16. Thoring A., Rudolph D., \& Vogl R. Digitalization of higher
education from a student's point of view//European Journal of Higher
Education IT. -2017.- Vol.47.-P.1-10. URL:
\url{https://www.semanticscholar.org/paper/Digitalization-of-Higher-Education-from-a-Student-\%E2\%80\%99-Thoring-Rudolph/0da7cba8d0eaeb81cbbfe792637d0983365e6d46}.-
Date of access: 06.01 2026.

17. Bleeker A., \& Crowder R. Selected online learning experiences in
the Caribbean during COVID-19//Studies and Perspectives Series-ECLAC.
United Nations. -2022.-No.105.\\
URL:\url{https://hdl.handle.net/11362/47742}. - Date of access: 06.01
2026.

18. Нахбаева Г. С., Нуртазина Р. А., \& Кайдарова А. С. Цифровая
трансформация вузов в Казахстане: результаты глубинных интервью с
экспертами// Адам әлемі. -2025.- \emph{Т.1}(103).- С.101-112. DOI
10.48010/aa.v103i1.724.

19. Мусина С. К. Применение цифровых инструментов в управлении высшими
учебными заведениями Казахстана// Управление образованием. -2024.-
Т.14(8--1). - С.278-285.

20. Бекманова Г.Т., Кантуреева М.А., Омарбекова А.С., Ергеш Б. Ж., \&
Закирова А.Б. Жоғары білім беруде жасанды интеллектіні қолдану және
әсерін бағалау// Academic Scientific Journal of Computer Science.
-2025.- Т.2.-№354.- C.61-73. DOI 10.32014/2025.2518-1726.344.

21. Министерство науки и высшего образования Республики Казахстан.
Межвузовский стандарт по применению искусственного интеллекта в высшем и
послевузовском образовании.-2024.-URL:
\href{https://enic-kazakhstan.edu.kz/files/1703569802/o-proekte-mezhvuzovskogo-standarta-po-primeneniyu-iskusstvennogo-intellekta.pdf}{https://enic-kazakhstan.edu.kz}. Дата обращения: 06.01.2026.

22. Үкімет Жасанды интеллектіні дамытудың 2024-2029 жылдарға арналған
тұжырымдамасын қабылдады.URL:
\href{https://primeminister.kz/news/ukimet-zhasandy-intellektini-damytudyn-2024-2029-zhyldarga-arnalgan-tuzhyrymdamasyn-kabyldady-28786}{https://primeminister.kz}. -Қаралған күні: 06.01.2026.

23. Inform.kz. ИИ в образовании: 95 вузов Казахстана интегрировали новые
дисциплины.-2025.-URL:\url{https://www.inform.kz/ru/ii-v-obrazovanii-95-vuzov-kazahstana-integrirovali-novie-distsiplini-f8be73}.
-- Дата обращения: 06.01 2026.

24. Alt University. Алматыда жоғары оқу орындары басшыларына арналған
«Жоғары білім беру жүйесінде жасанды интеллектіні пайдалану»
бағдарламасы бойынша оқу өтті. -2025.-URL:
\href{https://alt.edu.kz/zhanalyqtar/almatyda-zhoghary-oqu-oryndary-basshylaryna-arnalghan-zhoghary-bilim-beru-zhuejesinde-zhasandy-intellektini-pajdalanu-baghdarlamasy-bojynsha-oqu-oetti/}{https://alt.edu.kz} .-
Қаралған күні: 06.01.2026.

25. Korkyt Ata University. AI-Sana: Қазақстанда жасанды интеллект
болашағын қалыптастыратын инновациялық бағдарлама. -2025.-URL:
\url{https://korkyt.edu.kz/article/1806}.- Қаралған күні: 06.01.2026.

26. ҚСЗИ. Қазақстандықтар көбінесе цифрлық құралдарды оқу үрдісі мен
өзін-өзі дамыту үшін пайдаланады -- ҚСЗИ әлеуметтанулық сауалнамасы.
-2025.- URL:
\href{https://kisi.kz/kazakstandyktar-kobinese-chifrlyk-kuraldardy-oku-urdisi-men-ozin-ozi-damytu-ushin-pajdalanady-kszi-aleumettik-saualnamasy/}{https://kisi.kz} .-
Қаралған күні: 06.01.2026.

27. Heppi. AI and the Future of Universities. -2025.-URL:
https://www.hepi.ac.uk/reports/right-here-right-now-new-report-on-how-ai-is-transforming-higher-education/.-
Date of access: 06.01 2026.

28. Duke University. AI at Duke. URL: \url{https://ai.duke.edu/}.- Date
of access: 06.01 2026.

29. Contractor S., \& Reyes M. Generative AI in Higher Education:
Evidence from an Elite College// Higher Education \& Technology Review.
-2025.- Vol.12(1).- P.77--94.\ul{\hfill\break
}DOI 10.48550/arXiv.2508.00717.

30. Moldashev K., \& Sakhimbekov B. Determinants of teachers' adoption
of artificial intelligence: Evidence from Kazakhstan// Eurasian Journal
of Economic and Business Studies.- 2025.DOI 10.47703/ejebs.v69i4.565.
\end{refs}

\begin{center}
{\bfseries References}
\end{center}

\begin{refs}
1. George B., \& Wooden O. Managing the Strategic Transformation of
Higher Education through Artificial Intelligence//~Administrative
Sciences. -2023.-~Vol.13(9). -P.196. DOI 10.3390/admsci13090196.

2. Kane G. C. Digital maturity, not digital transformation// MIT Sloan
Management Review. -2017.- Vol.1(1). -P.1--15. URL:
\url{https://sloanreview.mit.edu/article/digital-maturity-not-digital-transformation}.
- Date of access: 06.01.2026.

3. Zakon.kz. Қazaқstan Respublikasynyң Zhoғary zhәne (nemese) zhoғary
oқu ornynan kejіngі bіlіm beru ұjymdarynyң cifrlyқ zhetіlu deңgejіn
ajқyndau zhөnіndegі әdіstemelіk ұsynymdardy bekіtu turaly// Қazaқstan
Respublikasy Ғylym zhәne zhoғary bіlіm ministrіnің mіndetіn atқarushynyң
2024 zhylғy 14 tamyzdaғy № 408 bұjryғy. -2024. -
URL: https://online.zakon.kz/Document/?doc\_id=34078184.-Қaralғan kүnі:
06.01.2026. {[}in Kazakh{]}

4. Dzhanegizova A., Huong Lan D., \& Chulanova Z. Cifrovaja
transformacija vysshego obrazovanija v Kazahstane: problemy i
reshenija// Statistika, uchet i audit. - 2025. -T.3(98). - S.
249-263.DOI 10.51579/1563-415.2025.-3.18.{[}in Russian{]}

5. Antonopoulou K., Begkos C., \& Zhu Z. Staying afloat amidst extreme
uncertainty: A case study of digital transformation in higher
education// Technological Forecasting \& Social Change.
-2023.-Vol.192:122603. DOI 10.1016/j.techfore.2023.122603.

6. Hess T., Matt C., Benlian A., \& Wiesböck F. Options for formulating
a digital transformation strategy. MIS Quarterly Executive. -2016.-
Vol.15(2). -P.123-139. URL:
\url{https://aisel.aisnet.org/cgi/viewcontent.cgi?article=1335&context=misqe}.
Date of access: 06.01.2026.

7. Shrivastava S. K., \& Shrivastava C. New reforms in Indian higher
education: Rethinking the education model//Indian Journal of Social
Science and Literature. -2024.-Vol.4(1). -P.42-48.\\
DOI 10.54105/ijssl. E1142.04010924.

8. Yesner R. The Future of Higher Education: Digital Transformation Is
Critical to Learner and Institution Success// IDC.-2020.-URL:\\
\url{https://www.ecampusnews.com/files/2020/10/IDC_The-Future-of-Higher-Education.pdf}
-Date of access: 06.01 2026.

9. Stokes P., Baker N., DeMillo R., Soares L., \& Yaeger L. The
transformation-ready higher education institution. Huron; American
Council of Education; Georgia Institute of Technology. -2019.- URL:
\url{https://share.google/wbxeAYP6ZYMYlnZSn}.- Date of access: 06.01
2026.

10. Rodríguez D., Arias P., \& Jimenez M. Use of virtual resources for
the development of spatial vision in high school students//Cadernos de
Educação Tecnologia e Sociedade. -2023.- Vol.16(3). DOI
10.14571/brajets. v16.n3.473-489.

11. Yureva O. V., Burganova L. A., Kukushkina O. Y., Myagkov G. P., \&
Syradoev D. V. Digital transformation and its risks in higher education:
Students' and teachers' attitude//Universal Journal of Educational
Research.-2020.-Vol.8(11B). -P.5965-5971. DOI 10.13189/ujer.2020.082232.

12. Wang S., Wang F., Zhu Z., Wang J., Tran T., \& Du Z. Artificial
intelligence in education: A systematic literature review//Expert
Systems with Applications. -2024.- Vol.252:124167.\\
DOI 10.1016/j.eswa.2024.124167.

13. Bond M., Khosravi H., De Laat M., et al. A meta systematic review of
artificial intelligence in higher education: A call for increased
ethics, collaboration, and rigour//International Journal of Educational
Technology in Higher Education\emph{.} -2024.-Vol.21:4. DOI
10.1186/s41239-023-00436-z.

14. Boussouf Z., Amrani H., Zerhouni Khal M., \& Daidai F. Artificial
intelligence in education: A systematic literature review// Data and
Metadata\emph{.} -2024.- Vol.3: 288. DOI 10.56294/dm2024288.

15. Tlili A. Can artificial intelligence (AI) help in computer science
education? A meta-analysis approach// Revista Española de Pedagogía.
-2024. DOI 10.22550/2174-0909.4172.

16. Thoring A., Rudolph D., \& Vogl R. Digitalization of higher
education from a student's point of view//European Journal of Higher
Education IT. -2017.- Vol.47.-P.1-10. URL:
\url{https://www.semanticscholar.org/paper/Digitalization-of-Higher-Education-from-a-Student-\%E2\%80\%99-Thoring-Rudolph/0da7cba8d0eaeb81cbbfe792637d0983365e6d46}.-
Date of access: 06.01 2026.

17. Bleeker A., \& Crowder R. Selected online learning experiences in
the Caribbean during COVID-19//Studies and Perspectives Series-ECLAC.
United Nations. -2022.-No.105.\\
URL:\url{https://hdl.handle.net/11362/47742}. - Date of access: 06.01
2026.

18. Nahbaeva G. S., Nurtazina R. A., \& Kajdarova A. S. Cifrovaja
transformacija vuzov v Kazahstane: rezul' taty glubinnyh
interv' ju s jekspertami// Adam әlemі. -2025.- T.1(103).
- S.101-112. DOI 10.48010/aa.v103i1.724. {[}in Russian{]}

19. Musina S. K. Primenenie cifrovyh instrumentov v upravlenii vysshimi
uchebnymi zavedenijami Kazahstana// Upravlenie obrazovaniem. -2024.-
T.14(8-1). - S.278-285. {[}in Russian{]}

20. Bekmanova G.T., Kantureeva M.A., Omarbekova A.S., Ergesh B. Zh., \&
Zakirova A.B. Zhoғary bіlіm berude zhasandy intellektіnі қoldanu zhәne
әserіn baғalau// Academic Scientific Journal of Computer Science.
-2025.- T.2.-№354.- C.61-73. DOI 10.32014/2025.2518-1726.344. {[}in
Russian{]}

21. Ministerstvo nauki i vysshego obrazovanija Respubliki Kazahstan.
Mezhvuzovskij standart po primeneniju iskusstvennogo intellekta v
vysshem i poslevuzovskom obrazovanii. -2024.-URL:
\href{https://enic-kazakhstan.edu.kz/files/1703569802/o-proekte-mezhvuzovskogo-standarta-po-primeneniyu-iskusstvennogo-intellekta.pdf}{https://enic-kazakhstan.edu.kz}.
-Data obrashhenija: 06.01.2026.{[}in Russian{]}

22. Үkіmet Zhasandy intellektіnі damytudyң 2024-2029 zhyldarғa arnalғan
tұzhyrymdamasyn қabyldady.URL:
\href{https://primeminister.kz/news/ukimet-zhasandy-intellektini-damytudyn-2024-2029-zhyldarga-arnalgan-tuzhyrymdamasyn-kabyldady-28786}{https://primeminister.kz}.
-Қaralғan kүnі: 06.01.2026.{[}in Kazakh{]}

23. Inform.kz. II v obrazovanii: 95 vuzov Kazahstana integrirovali novye
discipliny.-2025.-URL:
https://www.inform.kz/ru/ii-v-obrazovanii-95-vuzov-kazahstana-integrirovali-novie-distsiplini-f8be73.
- Data obrashhenija: 06.01 2026. {[}in Russian{]}

24. Alt University. Almatyda zhoғary oқu oryndary basshylaryna arnalғan
«Zhoғary bіlіm beru zhүjesіnde zhasandy intellektіnі pajdalanu»
baғdarlamasy bojynsha oқu өttі. -2025.-URL:
https://alt.edu.kz/zhanalyqtar/almatyda-zhoghary-oqu-oryndary-basshylaryna-arnalghan-zhoghary-bilim-beru-zhuejesinde-zhasandy-intellektini-pajdalanu-baghdarlamasy-bojynsha-oqu-oetti/.-
Қaralғan kүnі: 06.01.2026.{[}in Kazakh{]}

25. Korkyt Ata University. AI-Sana: Қazaқstanda zhasandy intellekt
bolashaғyn қalyptastyratyn innovacijalyқ baғdarlama. -2025.-URL:
https://korkyt.edu.kz/article/1806.- Қaralғan kүnі: 06.01.2026. {[}in
Kazakh{]}

26. ҚSZI. Қazaқstandyқtar kөbіnese cifrlyқ құraldardy oқu үrdіsі men
өzіn-өzі damytu үshіn pajdalanady -- ҚSZI әleumettanulyқ saualnamasy.
-2025.- URL:
\href{https://kisi.kz/kazakstandyktar-kobinese-chifrlyk-kuraldardy-oku-urdisi-men-ozin-ozi-damytu-ushin-pajdalanady-kszi-aleumettik-saualnamasy/}{https://kisi.kz} .-
Қaralғan kүnі: 06.01.2026. {[}in Kazakh{]}

27. Heppi. AI and the Future of Universities. -2025.-URL:
https://www.hepi.ac.uk/reports/right-here-right-now-new-report-on-how-ai-is-transforming-higher-education/.-
Date of access: 06.01 2026.

28. Duke University. AI at Duke. URL: \url{https://ai.duke.edu/}.- Date
of access: 06.01 2026.

29. Contractor S., \& Reyes M. Generative AI in Higher Education:
Evidence from an Elite College// Higher Education \& Technology Review.
-2025.- Vol.12(1).- P.77--94.\ul{\hfill\break
}DOI 10.48550/arXiv.2508.00717.

30. Moldashev K., \& Sakhimbekov B. Determinants of teachers' adoption
of artificial intelligence: Evidence from Kazakhstan// Eurasian Journal
of Economic and Business Studies.- 2025.DOI 10.47703/ejebs.v69i4.565.
\end{refs}

\begin{info}
\hspace{1em}\emph{{\bfseries Авторлар туралы мәліметтер}}

Торебекова З.Т. - аға ғылыми қызметкер, Алматы менеджмент
университеті, Алматы, Қазақстан, e-mail: zttemirkhan@gmail.com;

Бокаев Б.Н. - профессор, Қарағанды индустриялық университеті,
Теміртау, Қазақстан, e-mail: baurzhanbokayev@gmail.com;

Бактияров А.Д. - PhD докторанты, Л.Н. Гумилев атындағы Еуразиялық
ұлттық университет, Астана, Қазақстан, e-mail:
baktiyarov.asset@gmail.com;

Айжарыков Н.Б. - Зерттеу тобының мүшесі, Алматы менеджмент
университеті, Алматы, Қазақстан, e-mail: nblectures@gmail.com.

\hspace{1em}\emph{{\bfseries Information about the authors}}

Torebekova Z.-Senior Research Fellow, Almaty Management University,
Almaty, Kazakhstan, e-mail: zttemirkhan@gmail.com;

Bokayev B.- Professor, Karaganda Industrial University, Temirtau,
Kazakhstan, e-mail: baurzhanbokayev@gmail.com;

Baktiyarov A. - PhD Doctoral Student, L.N. Gumilyov Eurasian National
University, Astana, Kazakhstan, e-mail: baktiyarov.asset@gmail.com;

Aizharykov N. - Member of the Research Group, Almaty Management
University, Almaty, Kazakhstan, e-mail: nblectures@gmail.com.
\end{info}
